\uuid{WPHK}
\titre{Congruences}
\auteur{bodin ; croizat ; sacre}
\organisation{exo7}
\licence{CC BY-NC-SA 4.0}
\datecreate{2022-03-08}
\niveau{L1}
\module{Algèbre}
\chapitre{Arithmétique dans Z}
\sousChapitre{Autre}
\difficulte{2}
\theme{Congruences}
\exo7qcmid{hilisit.arithmetique.046}

\contenu{
\begin{question}
\qclassification{103.99}
 Soit $n$ un entier premier avec $3$. On peut alors affirmer :
    \begin{answers}
        \bad{$ 2n \equiv 1 \;[3]$}
        \bad{$2n \equiv -1\; [3]$}
        \good{$n^2 \equiv 1 \;[3]$}
        \bad{$n^2 \equiv -1 \;[3]$}
    \end{answers}
    \begin{explanations}
    Puisque $n$ n'est pas un multiple de $3$, on a soit $n \equiv 1 \;[3]$ (cas 1) soit $n \equiv 2 \equiv -1 \;[3]$ (cas 2). Dans le cas 1, on a $2n \equiv 2 \equiv -1 \;[3]$, et dans le cas 2 on a $2n \equiv -2 \equiv 1 \; [3]$.\\
    Dans les deux cas, on aura $n^2 \equiv 1 \; [3]$.
    \end{explanations}
\end{question}
}
